% !TeX root = ../main.tex
\clearpage
\section{Baseline recovery and peak detection in simulated air-quality data}
\label{sec:air_simulation}

We next consider a setting in which the observed series combines a drifting baseline with short-lived peaks. This experiment follows the peak-detection design in Section~4.2 of \citet{brantley2020baseline}, motivated by air-quality sensors whose slowly varying background can obscure transient pollution events. The objective is to estimate a low-quantile baseline and then identify peaks in the detrended observations, so the experiment examines both curve estimation and a downstream classification task. We compare Bayesian-PMALA with the horseshoe Gibbs sampler of \citet{onizuka2024fast} and the optimization benchmark labeled ADMM. Each quantile level is fitted separately.

\subsection{Data-generating mechanism and estimation target}
\label{subsec:air_design}

For each $n\in\{500,1000,2000,4000\}$, we generate $100$ series at observation indices $t_i=i$, $i=1,\ldots,n$, according to
\begin{equation}
 y_i=b_i+s_i+e_i,\qquad e_i\stackrel{\mathrm{iid}}{\sim}\mathcal{N}(0,0.25^2).
 \label{eq:air_observation}
\end{equation}
Here $b_i$ is the baseline drift and $s_i$ is the superposition of transient peaks. Both components are regenerated for each replication. For the baseline, let $u_i=(i-0.5)/n$ and draw $d=1+Z$, where $Z\sim\operatorname{Poisson}(n/100)$. Using the natural cubic spline basis returned by \texttt{splines::ns} with $d$ degrees of freedom and its default knot placement, set
\begin{equation}
 b_i=\sum_{\ell=1}^{d}a_\ell B_\ell(u_i),
 \qquad a_\ell\stackrel{\mathrm{iid}}{\sim}\operatorname{Exp}(1).
 \label{eq:air_baseline}
\end{equation}
The $1+$ offset in $d$ follows the supplied generator. The number of basis functions, and hence baseline complexity, varies across replications and increases on average with $n$.

Independently, draw the number of peaks $K\sim\operatorname{Binomial}(n,0.005)$. For $j=1,\ldots,K$, generate $U_j\sim\operatorname{Unif}(0,1)$, peak centers $c_j=\operatorname{round}\{1+(n-2)U_j\}$, widths $w_j\sim\operatorname{Unif}(2,12)$, and amplitudes $A_j\sim\mathcal{N}(20,4^2)$. The peak component is
\begin{equation}
 s_i=\sum_{j=1}^{K}\frac{A_j}{w_j}\,
 \phi\!\left(\frac{i-c_j}{w_j}\right),
 \label{eq:air_peaks}
\end{equation}
where $\phi$ is the standard normal density and an empty sum equals zero. Thus $A_j$ multiplies a Gaussian density and controls its area, rather than its height. The amplitudes are used without truncation. Peak widths are measured in observation-index units, whereas the spline is constructed on the normalized grid $u_i$. The expected peak counts are $2.5$, $5$, $10$, and $20$ for the four sample sizes, respectively. Peaks may overlap, and some small-sample replications contain no points above the signal threshold defined below. Since both the baseline complexity and expected peak count change with $n$, this is not a fixed-curve experiment in which only the observation grid becomes denser.

We fit the levels $\tau\in\{0.01,0.05,0.10\}$. Conditional on the generated baseline, the reference curve is the quantile of the peak-free process $y_i-s_i$:
\begin{equation}
 q_{\tau,i}^{0}=b_i+0.25\Phi^{-1}(\tau),
 \label{eq:air_target}
\end{equation}
where $\Phi$ is the standard normal distribution function. The methods receive the contaminated observations $y_i$, without access to $b_i$ or $s_i$. The true components are used only for evaluation. The target in \eqref{eq:air_target} differs from the conditional quantile $b_i+s_i+0.25\Phi^{-1}(\tau)$ of the contaminated series with its realized peaks held fixed. Low-quantile fitting is used to recover the background despite that contamination, rather than to track the peaks themselves.

The same generated series is used by all methods and all three quantiles for a given sample size and replication. The completed comparison therefore contains $1200$ fitted configurations per method, arising from $400$ generated datasets and three quantiles. Summary results use one completed campaign per method.

\subsection{Fitting procedures and numerical settings}
\label{subsec:air_fitting}

All methods use trend-filtering order $k=2$, corresponding to a third-difference penalty and piecewise-quadratic fitted trends. Bayesian-PMALA is the sampler of Section~\ref{subsec:individual} with the configuration used in Section~\ref{sec:simulations}: $\sigma\sim\mathrm{IG}(0.1,0.1)$, a chain started at $\theta=y$, $30000$ sweeps of which the first $10000$ are discarded, and the numerical settings of Appendix~\ref{app:pmala}. The only difference is the prior on $\lambda$.

The data model of this experiment, $y_i=b_i+s_i+e_i$, differs from $y_i=\theta_i+\varepsilon_i$ of Section~\ref{sec:simulations} in a way that the likelihood cannot resolve. The target is the baseline beneath the bursts, and the asymmetric Laplace likelihood at a low quantile level does not distinguish a burst from a rapid movement of the trend. Under the weak prior $\mathrm{Gamma}(1,10^{-2})$ the posterior accommodates the bursts in $\theta$: the effective penalty $\lambda\sigma$ settles well below the penalty that extended BIC selects for the same fits, and the posterior mean follows the peaks. The separation between a smooth background and transient excursions is a structural assumption about the data, and in a Bayesian analysis it enters through the prior. For series with bursts we therefore use
\begin{equation}
  \lambda\sim\mathrm{Gamma}\!\left(m,\;m/\lambda^\ast\right),\qquad m=n-k-1,
  \label{eq:air_prior}
\end{equation}
with the same sampler. The prior mean is the constant $\lambda^\ast$, which does not depend on $n$ or on the data. The shape $m$ gives the prior the same weight as the $m$ difference terms in the conditional update $\lambda\mid\theta\sim\mathrm{Gamma}(a+m,\,b+\lVert D^{(k+1)}\theta\rVert_1)$, so the penalty remains random, with posterior coefficient of variation near $(2m)^{-1/2}$. We set $\lambda^\ast=5000$ at the extreme level $\tau=0.01$ and $\lambda^\ast=30000$ at the inner levels $\tau\in\{0.05,0.10\}$. The extreme level requires the smaller value because a penalty that is too strong for the few observations below the curve inflates $\sigma$, which raises the effective penalty $\lambda\sigma$ further until the trend move stops accepting. Each constant is the largest value in a screening grid for which no chain stopped moving at any $n\in\{500,\ldots,4000\}$; the screening used replications $1$--$3$ at every sample size and replications $1$--$20$ at $n=1000$. Because $\lambda$ has the units of $1/y$, the constants refer to the scale of this design, whose noise standard deviation is $0.25$; Section~\ref{sec:app:spod} applies the prior to data on another scale.

The horseshoe competitor uses the Gibbs sampler of \citet{onizuka2024fast}, retaining its difference-operator scaling, initialization, conditional updates, and numerical bounds on scale parameters. The campaign does not run the supplied code itself but a re-implementation of those conditionals, in which the trend is drawn through a sparse Cholesky factorization of its banded precision matrix and the observation-level latent variables are drawn in a vectorized manner; the re-implementation was checked against the supplied code before use. The completed campaign uses $40000$ iterations with $10000$ burn-in iterations, increased from an $8000$-iteration pilot after longer-chain checks revealed sensitivity to run length. Its point estimate averages all $30000$ retained draws. Both Bayesian procedures are summarized by posterior means.

The ADMM-labeled benchmark selects a penalty from
\[
 \Lambda=\{\exp(-4+j/2):j=0,\ldots,24\}
\]
by the implemented extended BIC criterion. For a candidate fit, let $L_\tau=\sum_i\rho_\tau(y_i-\widehat\theta_{\tau,i})$ and let $\nu$ count third differences whose absolute values exceed $10^{-4}$. The selection score is
\begin{equation}
 \operatorname{eBIC}=\frac{2L_\tau}{\min(\tau,1-\tau)}
       +\nu\log n+2\log\binom{n-3}{\nu}.
 \label{eq:air_ebic}
\end{equation}
The fitting code minimizes the sum of check losses plus $\lambda\|D^{(3)}\theta\|_1$. Its configured window size is $4000$, so every series here is solved as a single sparse linear program using \texttt{PIQP}; the multi-window ADMM branch is not invoked. The solver checks a feasible dual bound against a relative objective tolerance of $10^{-5}$. We retain the source label ADMM in the tables to identify this benchmark, but it denotes the independently fitted quantile trend-filtering estimator in this campaign. It is not the jointly constrained, noncrossing \texttt{detrend\_eBIC} comparison reported in the original paper.

\subsection{Evaluation criteria}
\label{subsec:air_metrics}

For a fitted curve $\widehat\theta_{\tau,i}^{(r)}$ in replication $r$, we compute
\begin{align}
 \operatorname{MSE}^{(r)}_\tau
  &=\frac{1}{n}\sum_{i=1}^n
    \bigl(\widehat\theta_{\tau,i}^{(r)}-q_{\tau,i}^{0,(r)}\bigr)^2,
    \label{eq:air_mse}\\
 \operatorname{CL}^{(r)}_\tau
  &=\frac{1}{n}\sum_{i=1}^n
    \rho_\tau\!\bigl(y_i^{(r)}-\widehat\theta_{\tau,i}^{(r)}\bigr).
    \label{eq:air_check_loss}
\end{align}
MSE measures recovery of the peak-free quantile curve. Check loss (CL) measures in-sample fit to the contaminated observations. Consequently, a lower CL need not imply more accurate baseline recovery: a procedure can reduce the observed-data loss while absorbing transient structure that should remain after detrending. Comparisons of CL are made within a quantile level.

For peak detection, define the true and estimated signal indicators by
\begin{equation}
 z_i^{(r)}=\mathbbm{1}\{s_i^{(r)}>0.5\},\qquad
 \widehat z_{\tau,h,i}^{(r)}=
 \mathbbm{1}\{y_i^{(r)}-\widehat\theta_{\tau,i}^{(r)}>h\},
 \label{eq:air_classification}
\end{equation}
using $h\in\{0.5,0.75,1,1.2\}$. These are the thresholds used in the completed campaigns; the design is adapted from Section~4.2 of \citet{brantley2020baseline} without asserting that all four threshold values reproduce its original comparison. We subtract the fitted low-quantile curve directly, without an additional Gaussian quantile-offset correction. Thus the quantile and classification threshold jointly determine the detector.

Because signal points are relatively rare, we summarize classification by class-averaged accuracy,
\begin{equation}
 \operatorname{CAA}^{(r)}_{\tau,h}=
 \frac{1}{2}\left\{
 \frac{\sum_i z_i^{(r)}\widehat z_{\tau,h,i}^{(r)}}{\sum_i z_i^{(r)}}+
 \frac{\sum_i(1-z_i^{(r)})(1-\widehat z_{\tau,h,i}^{(r)})}
      {\sum_i(1-z_i^{(r)})}\right\}.
 \label{eq:air_caa}
\end{equation}
CAA weights sensitivity and specificity equally. A perfect classifier attains one; predicting no signal everywhere gives $0.5$ when both classes are present. We first calculate each metric within a replication and then take its arithmetic mean across replications. Undefined values are omitted only from the affected metric. In particular, some $n=500$ datasets contain no points with $s_i>0.5$, and CAA is undefined for them. No fit is excluded from the averages because of a mixing diagnostic. The main tables report replication means; the standard errors across replications, and the paired differences between methods that the shared datasets permit, are given in Table~\ref{tab:air_mse_se}.

\subsection{Results and interpretation}
\label{subsec:air_results}

Table~\ref{tab:air_n1000} reports the results for $n=1000$. The remaining sample sizes are given in Appendix~\ref{app:air_simulation} (Tables~\ref{tab:air_n500}, \ref{tab:air_n2000}, and \ref{tab:air_n4000}), together with standard errors and paired comparisons on the shared datasets (Table~\ref{tab:air_mse_se}), finite-chain and interval diagnostics (Table~\ref{tab:air_diagnostics}), and penalty levels (Table~\ref{tab:air_penalty}).

% AIR-TABLE-1000-BEGIN
\begin{table}[H]
\centering
\small
\setlength{\tabcolsep}{4pt}
\caption{Simulated air-quality results for $n=1000$. Entries are means across replications. CL denotes check loss; HS denotes the horseshoe sampler.}
\label{tab:air_n1000}
\begin{tabular}{@{}rlrrrrrr@{}}
\toprule
 & & & & \multicolumn{4}{c}{CAA at threshold $h$} \\
\cmidrule(lr){5-8}
$\tau$ & Method & CL & MSE & $0.5$ & $0.75$ & $1$ & $1.2$ \\
\midrule
0.01 & Bayesian-PMALA & 0.005652 & 0.045406 & 0.6975 & 0.7805 & 0.7641 & 0.7137 \\
0.01 & Onizuka (HS) & 0.005227 & 0.100689 & 0.5872 & 0.6272 & 0.6136 & 0.5909 \\
0.01 & ADMM & 0.006652 & 0.018925 & 0.6950 & 0.8364 & 0.8701 & 0.8125 \\
\addlinespace[3pt]
0.05 & Bayesian-PMALA & 0.028083 & 0.012385 & 0.7936 & 0.8605 & 0.8076 & 0.7336 \\
0.05 & Onizuka (HS) & 0.023669 & 0.099941 & 0.5736 & 0.5779 & 0.5517 & 0.5347 \\
0.05 & ADMM & 0.028930 & 0.009350 & 0.8014 & 0.8849 & 0.8400 & 0.7610 \\
\addlinespace[3pt]
0.10 & Bayesian-PMALA & 0.049294 & 0.011511 & 0.8358 & 0.8495 & 0.7688 & 0.7025 \\
0.10 & Onizuka (HS) & 0.041807 & 0.100682 & 0.5705 & 0.5570 & 0.5324 & 0.5222 \\
0.10 & ADMM & 0.050643 & 0.009270 & 0.8478 & 0.8761 & 0.7986 & 0.7238 \\
\bottomrule
\end{tabular}
\end{table}
% AIR-TABLE-1000-END

Bayesian-PMALA recovers the peak-free baseline far more accurately than the locally adaptive horseshoe sampler, in every configuration and in most replications, and it detects the peaks correspondingly better. The local shrinkage scales that made the horseshoe competitive on the varying-smoothness signal of Section~\ref{sec:simulations} allow its fitted curve to bend toward each peak; this lowers the check loss on the contaminated observations but raises the error against the peak-free target, and the detector, which subtracts the fitted curve, then removes part of the signal it is meant to find. A single global penalty cannot bend locally in this way, and in this design that rigidity is the appropriate behaviour. The ADMM benchmark, designed for this problem, remains the most accurate method, and Bayesian-PMALA comes close to it in both baseline recovery and detection while providing posterior inference on the baseline, the penalty and the likelihood scale. Its penalty is set by a prior fixed before the data are seen, whereas the ADMM penalty is chosen by extended BIC over a fixed grid and attains the largest grid value in a substantial fraction of the fits at $\tau=0.05$ and $0.10$ (Table~\ref{tab:air_penalty}). The acceptance rate of the trend move is close to its target of $0.5$ in every configuration, and the pointwise $95\%$ bands of both samplers cover the peak-free target less often than nominal, since the posterior conditions on the contaminated series (Table~\ref{tab:air_diagnostics}).
